% Antecedentes recuperados. Fragmento autónomo en sus definiciones;
% requiere amsmath, amssymb, amsthm, booktabs y longtable del documento principal.
\section{Réalisation finie du transport avec mémoire de phase}
\label{iii:hist:antecedentes}

Les deux développements suivants réunissent des mécanismes complémentaires de la
construction : un transfert fini sur des positions APP, des modes et une phase,
et une collection de lecteurs cinématiques et constitutifs sur deux canaux
angulaires. Chaque réalisation conserve son espace d’états, ses opérations
et ses paramètres. L’opérateur fini a le domaine défini
ci-dessous ; le TPK enrichi du corps principal conserve en outre ses
propres coordonnées et prolongements. Les lecteurs réciproques sont présentés
comme des réalisations auxiliaires avec des formules et un domaine déclarés.

\subsection{Transfert fini sur les positions, les modes et la phase}
\label{iii:hist:transferencia}

Soit $\operatorname{dr}_9(n)$ le représentant dans $\{1,\ldots,9\}$ de
$n\bmod9$, avec $9$ comme représentant du résidu nul. Aux positions
$i,j\in\{0,\ldots,8\}$, on considère les deux tables
\[
 A_+(i,j)=\operatorname{dr}_9(i+j+2),\qquad
 A_\times(i,j)=\operatorname{dr}_9((i+1)(j+1)).
\]
Les entrées évaluées ici sont des entiers positifs. Ainsi, la convention
historique alternative $\operatorname{dr}_9(0)=0$ n’intervient pas dans ces tables.
L’état fini est
\[
 \Omega=(\mathbb Z/9\mathbb Z)^2\times\{+,\times\}
 \times\{N,E,S,O\}\times\mathbb Z/3\mathbb Z,
 \qquad |\Omega|=1944.
\]
Nous écrivons $x=(i,j,m,d,\tau)$. La transition conserve l’ordre suivant :
on choisit $d'$, on déplace la position, on lit APP avec le mode antérieur,
on réduit la lecture modulo $3$, on met à jour le mode et on avance la phase.
Avec
\[
 \delta_N=(-1,0),\quad\delta_E=(0,1),\quad
 \delta_S=(1,0),\quad\delta_O=(0,-1),
\]
les opérations sont
\begin{align*}
 (i',j')&=(i,j)+\delta_{d'}\pmod9,\\
 a&=A_m(i',j')\pmod3,\\
 m'&=\begin{cases}+,&a=1,\\\times,&a=2,\\m,&a=0,\end{cases}
 &\tau'&=\tau+1\pmod3.
\end{align*}
Ainsi, $F_{d'}(x)=(i',j',m',d',\tau')$. Le coût de rotation et de changement de mode est
\[
 c_\lambda(x,d')=\mathbf1_{d'\ne d}+\lambda\mathbf1_{m'\ne m},
 \qquad \beta>0,\quad\lambda\geq0,
\]
et l’opérateur agit sur la base d’états par
\begin{equation}
 T_{\beta,\lambda}e_x
 =\sum_{d'\in\{N,E,S,O\}}
 e^{-\beta c_\lambda(x,d')}e_{F_{d'}(x)}.
 \label{iii:hist:T}
\end{equation}
Les paramètres $\beta,\lambda$ sont des paramètres déclarés de ce modèle
fini. La table de référence utilise $(\beta,\lambda)=(1,1)$ ; le cas
$\lambda=0$ appartient aux tests de sensibilité.

\begin{proposition}[Équivariance de l’antécédent fini]
\label{iii:hist:equivarianza}
L’action
\[
 \Phi_{u,v,w}(i,j,m,d,\tau)
 =(i+3u,j+3v,m,d,\tau+w),\qquad
 (u,v,w)\in(\mathbb Z/3\mathbb Z)^3,
\]
commute avec $T_{\beta,\lambda}$.
\end{proposition}
\begin{proof}
On a $\operatorname{dr}_9(n)\equiv n\pmod3$. Par conséquent,
$A_+(i,j)\equiv i+j+2$ et
$A_\times(i,j)\equiv(i+1)(j+1)\pmod3$.
La translation par des multiples de trois préserve les deux lectures réduites.
Elle commute également avec le déplacement choisi et préserve le mode antérieur,
le mode postérieur et les deux directions. Elle conserve par conséquent chaque
indicateur du coût. Enfin, ajouter $w$ à la phase commute avec l’ajout d’une
unité. Chaque transition et son poids sont transportés vers la transition et le poids
correspondants ; la somme dans \eqref{iii:hist:T} conserve cette égalité.
\end{proof}

\begin{theorem}[Décomposition complète par caractères]
\label{iii:hist:descomposicion}
Pour $\omega=e^{2\pi i/3}$, l’espace $\mathbb C^\Omega$ se décompose
unitairement en vingt-sept sous-espaces de dimension $72$. Ses blocs
satisfont
\[
 T(k_a,k_b,k_t)=\omega^{k_t}T(k_a,k_b,0),
 \qquad k_a,k_b,k_t\in\{0,1,2\}.
\]
\end{theorem}
\begin{proof}
Nous écrivons $i=3a+r$, $j=3b+s$, avec $a,b,r,s\in\{0,1,2\}$.
Chaque fibre fine $u=(r,s,m,d)$ a $27$ états $(a,b,\tau)$.
Nous définissons l’isométrie $U_k:\mathbb C^{72}\to\mathbb C^\Omega$ par
\[
 U_ke_u=\frac1{\sqrt{27}}
 \sum_{a,b,\tau=0}^2
 \omega^{-(k_aa+k_bb+k_t\tau)}e_{(3a+r,3b+s,m,d,\tau)}.
\]
L’orthogonalité des caractères prouve $U_k^*U_\ell=\delta_{k\ell}I$ et
$\sum_kU_kU_k^*=I$. Si une transition franchit des frontières fines avec des accroissements
grossiers $\Delta a,\Delta b$, son coefficient dans $U_k^*TU_k$ est
\[
 e^{-\beta c_\lambda}\omega^{k_a\Delta a+k_b\Delta b+k_t}.
\]
Cela démontre $TU_k=U_kT(k)$ et la décomposition de l’opérateur. L’accroissement
temporel vaut toujours un, de sorte que le dernier facteur peut être extrait de
tout le bloc. Lorsque l’on élève au cube, ce facteur temporel disparaît, bien que
les caractères spatiaux conservent leur fonction.
\end{proof}

La décomposition précédente repose sur l’équivariance et l’orthogonalité
des caractères. La non-négativité du bloc trivial garantit en outre une
valeur propre réelle non négative égale au rayon spectral. Cette propriété
s’applique à l’espace complet, en conservant également ses états transitoires.

\subsection{Décomposition spectrale et sensibilité paramétrique}
\label{iii:hist:espectros}

Soient $\rho_1\geq\rho_2$ les deux plus grands modules, comptés avec
multiplicité, et $g=\rho_1-\rho_2$. Le tableau suivant contient les
neuf blocs spatiaux ; le théorème précédent détermine les vingt-sept
blocs temporels complets : pour chaque ligne et chaque $t=0,1,2$, la valeur propre
dominante est $\omega^t\mu_1(k_a,k_b,0)$, tandis que $\rho_1,\rho_2,g$
restent identiques. Cette forme conserve explicitement les trois phases.
Dans ce premier tableau, on fixe $\beta=\lambda=1$.

\begin{center}
\small
\begin{tabular}{ccrrrr}
\toprule
$k_a$&$k_b$&$\mu_1(k_a,k_b,0)$&$\rho_1$&$\rho_2$&$g$\\
\midrule
0&0&1.730888896&1.730888896&1.083076285&0.647812612\\
0&1&1.459962951&1.459962951&1.116843785&0.343119167\\
0&2&1.459962951&1.459962951&1.116843785&0.343119167\\
1&0&1.459962951&1.459962951&1.116843785&0.343119167\\
1&1&-1.508515719&1.508515719&1.001353790&0.507161929\\
1&2&-1.514010372&1.514010372&0.971977167&0.542033205\\
2&0&1.459962951&1.459962951&1.116843785&0.343119167\\
2&1&-1.514010372&1.514010372&0.971977167&0.542033205\\
2&2&-1.508515719&1.508515719&1.001353790&0.507161929\\
\bottomrule
\end{tabular}
\end{center}

Les signes des valeurs propres dominantes et leurs phases font partie de la
sortie : les quatre secteurs avec $k_a,k_b\ne0$ ont une valeur dominante réelle négative
pour $k_t=0$. Son module conserve le taux spectral et son signe conserve une
information supplémentaire qui intervient lors de la composition des itérations.

Les cinq applications complètes de sensibilité, avec les lignes $k_a$ et les colonnes
$k_b$, sont les suivantes ; toutes correspondent à $k_t=0$ :
\begin{align*}
\rho(1,1)&=\begin{pmatrix}
1.730888896&1.459962951&1.459962951\\
1.459962951&1.508515719&1.514010372\\
1.459962951&1.514010372&1.508515719
\end{pmatrix},\\
\rho(1,0)&=\begin{pmatrix}
2.103638324&1.685504327&1.685504327\\
1.685504327&1.853245966&1.853245966\\
1.685504327&1.853245966&1.853245966
\end{pmatrix},\\
\rho(1,2)&=\begin{pmatrix}
1.656915306&1.433869155&1.433869155\\
1.433869155&1.423885161&1.432116157\\
1.433869155&1.432116157&1.423885161
\end{pmatrix},\\
\rho(0.5,1)&=\begin{pmatrix}
2.480049445&2.124801655&2.124801655\\
2.124801655&2.281861809&2.286772914\\
2.124801655&2.286772914&2.281861809
\end{pmatrix},\\
\rho(2,1)&=\begin{pmatrix}
1.226049377&1.131390710&1.131390710\\
1.131390710&1.028041508&1.028201900\\
1.131390710&1.028201900&1.028041508
\end{pmatrix}.
\end{align*}
Ici, $\rho(\beta,\lambda)$ désigne l’application des rayons, non l’opérateur de
transfert. L’échange $i\leftrightarrow j$ accompagné de l’échange
correspondant des directions préserve APP, les transitions et le coût ;
ces applications sont donc symétriques. La variation des paramètres modifie les rayons tout en maintenant
la décomposition par caractères, qui dépend de la symétrie des
opérations. Les tableaux montrent des approximations numériques de spectres
finis. Les identités de transition, d’équivariance et de factorisation
reposent sur les preuves précédentes ; leurs conséquences numériques sont
conservées avec les données complètes des vingt-sept composantes.

\subsection{Lecteurs réciproques et tests de sensibilité}
\label{iii:hist:n6d}

Pour reproduire la réalisation angulaire antécédente, on conserve les mesures
numériques décimales en degrés
\[
 a_{\mathrm h}=7.297352569283802,\qquad
 c_{\mathrm h}=2.123738338968646.
\]
Les angles et la coordonnée adimensionnelle correspondante sont distingués par
\[
 A_{\mathrm h}=a_{\mathrm h}{}^\circ,\qquad
 C_{\mathrm h}=c_{\mathrm h}{}^\circ,\qquad
 \alpha_{\mathrm h}=a_{\mathrm h}/1000.
\]
Ces données permettent de confronter l’antécédent ; les constantes du développement
en vigueur conservent leur génération antérieure APP--TRIT--TPK et leur représentation
ultérieure. La reproduction historique ne remplace pas cette génération par les
décimales écrites ici. La conversion angulaire et les canaux utilisés sont
\[
 x_{\mathrm h}=\frac{\pi a_{\mathrm h}}{180},\qquad
 y_{\mathrm h}=\frac{\pi c_{\mathrm h}}{180},\qquad
 q_\pm=e^{-(x_{\mathrm h}\pm y_{\mathrm h})},
 \qquad0<q_+<q_-<1.
\]
Le premier lecteur exprime
\[
 \epsilon_{90}=q_-^{90}-q_+^{90},\qquad
 c_\pm=\frac{c_{\mathrm{ref}}}{1\pm\epsilon_{90}}.
\]
Ici, $c_{\mathrm{ref}}$ ne fait que normaliser le test cinématique. Pour reproduire
ses unités, la même valeur de référence historique a été utilisée,
$299792458\,\mathrm{m\,s^{-1}}$.

\begin{proposition}[Moyenne harmonique et produit de deux feuillets]
\label{iii:hist:lectores}
Si $|\epsilon_{90}|<1$, la moyenne harmonique de $c_+$ et $c_-$ est
$c_{\mathrm{ref}}$. Dans le lecteur constitutif historique
\[
 \mu_{\mathrm h}^{\pm}=e^{\pm2\alpha_{\mathrm h}},\qquad
 \varepsilon_{\mathrm h}^{\pm}=e^{\mp2\alpha_{\mathrm h}},
\]
on a $\mu_{\mathrm h}^{\pm}\varepsilon_{\mathrm h}^{\pm}=1$ et
$Z_+/Z_-=e^{4\alpha_{\mathrm h}}$.
\end{proposition}
\begin{proof}
La somme $c_+^{-1}+c_-^{-1}=2/c_{\mathrm{ref}}$ élimine le terme orienté.
Dans chaque feuillet, la somme des deux exposants constitutifs est nulle.
D’autre part, $Z_\pm=\sqrt{\mu_{\mathrm h}^{\pm}/
\varepsilon_{\mathrm h}^{\pm}}=e^{\pm2\alpha_{\mathrm h}}$.
Le quotient des impédances conserve donc l’orientation qu’élimine
le produit constitutif.
\end{proof}

Le deuxième lecteur historique provient de séries géométriques convergentes :
\begin{align*}
 S_a(q)&=\frac{q^a}{1-q^{3a}}
       =\sum_{j\geq0}q^{a(1+3j)},\\
 \delta_{a,b}(q)&=S_a(q)-S_b(q),\\
 R_{a,b}&=\frac{\delta_{a,b}(q_-)-\delta_{a,b}(q_+)}
 {q_-^a-q_+^a}.
\end{align*}
La convergence est valable pour $|q|<1$ ; le dénominateur du dernier quotient est
positif pour les canaux et les entiers positifs considérés. Les lecteurs
auxiliaires historiques sont
\[
 \zeta_{12}=1-R_{90,120}^2,\qquad
 \Delta_{4,\mathrm h}=\frac{x_{\mathrm h}-y_{\mathrm h}}{810}
 \left(1-\frac{7\pi}{12\cdot729}\right),\qquad
 \xi_{\mathrm h}=1-\frac{27}{160}\Delta_{4,\mathrm h}.
\]
L’indice $\mathrm h$ maintient explicite cette définition historique de la
torsion, en évitant de l’identifier par le symbole à un autre lecteur en vigueur.

L’évaluation à partir des chaînes décimales déclarées donne
\begin{align*}
 q_+&=0.848377942576404356\ldots,&
 q_-&=0.913660151150557274\ldots,\\
 \epsilon_{90}&=0.000295169439289827959\ldots,&
 R_{90,120}&=0.933314535237713061\ldots,\\
 \zeta_{12}&=0.128923978314011665\ldots,&
 \xi_{\mathrm h}&=0.999981235497802379\ldots.
\end{align*}
Le tableau complet des perturbations discrètes est
\begin{center}
\begin{tabular}{rrr}
\toprule
$a$&$b$&$R_{a,b}$\\\midrule
90&120&0.933314535238\\
90&119&0.927013603649\\
90&121&0.939071551912\\
89&120&0.939066204789\\
91&120&0.927019444194\\
30&40&0.569215041549\\
60&80&0.834174862952\\
\bottomrule
\end{tabular}
\end{center}
La sensibilité angulaire distingue les deux coordonnées de la paire :
\begin{center}
\begin{tabular}{lrrr}
\toprule
Variable&Variation relative&$R_{90,120}$&Variation en ppm\\\midrule
$A_{\mathrm h}$&$-0.001$&0.933059250359&$-273.525022305$\\
$A_{\mathrm h}$&$+0.001$&0.933568846564&$+272.481908792$\\
$C_{\mathrm h}$&$-0.001$&0.933388163678&$+78.889202964$\\
$C_{\mathrm h}$&$+0.001$&0.933240822366&$-78.979667757$\\
\bottomrule
\end{tabular}
\end{center}

Les évaluations ont été réalisées avec $80$ et $120$ chiffres de précision à partir
des chaînes décimales déclarées. La précision de l’évaluation conserve
la spécification finie de ces données angulaires.

Enfin, ces lecteurs maintiennent leur place dans la généalogie documentaire.
L’opérateur constitutif en vigueur
$V=(I-T^{90})(I-T^{120})^{-1}$ évalue
$(1-q_\pm^{90})/(1-q_\pm^{120})$ dans ses deux feuillets. La fonctionnelle orientée
en vigueur utilise $12S_{90}-S_{120}$. La réalisation antécédente conserve le lecteur
$S_{90}-S_{120}$ et sa normalisation par $q_-^{90}-q_+^{90}$.
La relation entre eux s’exprime en maintenant chaque opération et ses poids,
au lieu de les remplacer par un nombre unique ou un même symbole.
